\section{Objetivos del proyecto}

\subsection{Objetivo general}

De acuerdo con el Plan de Implementación aprobado, el objetivo principal es establecer la hoja de ruta estratégica y operativa para la adopción de herramientas digitales que fortalezcan la gestión comercial, la fidelización de clientes y la presencia internacional de Machu Picchu Exclusive Tours E.I.R.L. El proyecto se centra en la implementación de un CRM integral en la nube y en la ejecución de una estrategia de marketing digital multicanal que permita digitalizar la gestión de clientes, mejorar el seguimiento postventa y aumentar la captación en los mercados internacionales (norteamericano y europeo).

En el marco de este informe, dicho objetivo se materializa mediante la implementación de Odoo como plataforma central (CRM y módulos asociados) sobre la cual se articula el ecosistema digital de la empresa.

\subsection{Objetivos específicos}

\begin{itemize}
    \item[a.] \textbf{Implementar un CRM integral en la nube}: desplegar y poner en operación una plataforma CRM basada en Odoo que centralice la base de datos de clientes y el seguimiento de oportunidades comerciales, reemplazando la gestión dispersa en Excel, WhatsApp y registros físicos.
    \item[b.] \textbf{Digitalizar y automatizar el ciclo comercial}: integrar el CRM con los principales canales digitales (sitio web, redes sociales y WhatsApp Business) y habilitar campañas de marketing digital segmentadas que potencien la captación y fidelización de clientes internacionales.
    \item[c.] \textbf{Fortalecer la inteligencia de negocios}: definir y configurar indicadores clave de desempeño (KPIs) y tableros de control que permitan medir ventas, recurrencia de clientes, tiempos de respuesta y efectividad de las campañas digitales.
    \item[d.] \textbf{Desarrollar capacidades digitales del equipo}: ejecutar un plan de capacitación y gestión del cambio que asegure la adopción sostenible del CRM y de las herramientas de marketing digital por parte de la gerencia y el personal operativo.
    \item[e.] \textbf{Mejorar la competitividad y eficiencia operativa}: lograr, en el periodo de seguimiento establecido en el Plan de Implementación, incrementos en las ventas digitales y reducciones en los tiempos de respuesta y errores administrativos, alineados con las metas cuantitativas definidas en dicho plan.
\end{itemize}

